\section{Quadratic structure and a cubic multiplicity bound}\label{ranks}
For a torsion-free module $M$ of rank $r$ over a normal
domain, its \emph{reflexive determinant} is
\[
 \det M=\left(\bigwedge^r M\right)^{**},
 \qquad M^*=\Hom_R(M,R).
\]
If $R$ is local and $M$ has finite projective dimension, then
$\det M\simeq R$.  Indeed, choose a finite free resolution of $M$.
On the open set where $M$ is free, its determinant is the alternating
tensor product of the determinants of the free modules in this
resolution, hence is trivial.  That open set contains every point
of height at most one, since a torsion-free module over a height-one
localization of a normal domain is free over a discrete valuation ring.
An isomorphism of reflexive modules on this open set extends uniquely
across its complement, which has codimension at least two.

\subsection{A determinant refinement of quadratic splitting}

\begin{theorem}[Quadratic structure]\label{quadratic-order}
Let $R$ be a normal local domain satisfying $(R_2)$ and $(S_3)$, and
$R\subseteq S$ a torsion-free module-finite extension of rank two with
$\det S\simeq R$. Then $S^{**}=R\oplus R\theta$ is a quadratic algebra and
$S=R\oplus I\theta$ for an ideal with $I^{**}=R$. In particular the extension splits. Finite projective dimension of $S$ suffices for the determinant hypothesis.
\end{theorem}
\begin{proof}
Put $K=\operatorname{Frac}R$ and regard $S$ as a submodule
of its two-dimensional generic algebra $S_K=S\otimes_RK$.
Inside $S_K$ its reflexive hull is
\begin{equation}\label{quadratic-reflexive-intersection}
 T=S^{**}=\bigcap_{\operatorname{ht}\mathfrak p=1}S_{\mathfrak p}.
\end{equation}
Each member of this intersection is an algebra, so $T$ is an algebra
as well.  It is module-finite over $R$.  Consequently its elements
are integral over $R$, and normality gives $T\cap K=R$.
To check torsion-freeness of $T/R$, suppose $at\in R$ for a nonzero
$a\in R$ and $t\in T$. Then $t\in K\cap T=R$.

We next show that $T/R$ is reflexive.  At a prime of height at most
two, the base is regular and the reflexive module $T_{\mathfrak p}$
is free. If
$1_T\in\mathfrak pT_{\mathfrak p}$, multiplication would imply
$T_{\mathfrak p}=\mathfrak pT_{\mathfrak p}$, contradicting
Nakayama's lemma. Thus the unit is primitive and $(T/R)_{\mathfrak p}$ is free at these
primes.  At a prime of height at least three, $(S_3)$ gives
$\operatorname{depth}R_{\mathfrak p}\geq3$, while reflexivity of
$T$ gives $\operatorname{depth}T_{\mathfrak p}\geq2$.
The depth lemma applied to
\[
 0\longrightarrow R_{\mathfrak p}\longrightarrow T_{\mathfrak p}
   \longrightarrow(T/R)_{\mathfrak p}\longrightarrow0
\]
gives depth at least two for the quotient.  Hence $T/R$ is
torsion-free and satisfies $(S_2)$, which is the reflexivity
criterion over a normal domain.  This is precisely the quotient
step of \cite[Lemma~2.2]{HM}.

The map $t\mapsto1_T\wedge t$ induces a homomorphism
\[
 T/R\longrightarrow\det T=\det S.
\]
It is an isomorphism at every height-one prime, since there the
unit extends to a basis of the free rank-two module $T_{\mathfrak p}$.
Both sides are reflexive, so it is an isomorphism everywhere.
By the determinant hypothesis $T/R\simeq R$.  Lift a generator
of this quotient to $\theta\in T$.  Then
\begin{equation}\label{quadratic-order-description}
 T=R\oplus R\theta,
 \qquad \theta^2=\alpha+\beta\theta\quad(\alpha,\beta\in R).
\end{equation}
No division by two, separability assumption, or field hypothesis
on $S_K$ has been used.

For an element $a+b\theta\in S$, the scalar $a$ is already in $S$.
Subtracting it shows that
\[
 S=R\oplus I\theta,
 \qquad I=\{b\in R:b\theta\in S\}.
\]
The subset $I$ is an ideal.  Equation~\eqref{quadratic-reflexive-intersection}
gives $S_{\mathfrak p}=T_{\mathfrak p}$ at height one, so
$I_{\mathfrak p}=R_{\mathfrak p}$ there and $I^{**}=R$ as fractional
ideals.  Projection onto the scalar coordinate is the required
$R$-linear retraction.

The description also identifies the defect and the conductor:
\begin{equation}\label{quadratic-conductor}
 T/S\simeq R/I,
 \qquad (S:T)=\{t\in T:tT\subseteq S\}=IT.
\end{equation}
Indeed, writing $t=a+b\theta$, the condition $t\in S$ is $b\in I$,
and the condition $t\theta\in S$ is $a+\beta b\in I$.
Together they say $a,b\in I$, which is exactly $t\in IT$.
The finite-projective-dimension consequence follows from the
determinant observation at the beginning of this section.
\end{proof}

\subsection{Cubic algebras and Buchsbaum--Rim multiplicity}
In rank three the quotient by the unit has rank two, but need not
be free.  The proof measures its failure to be free by the length
of an Ext module.  The multiplicity input is the classical
parameter-module formula of Buchsbaum--Rim and its treatment in
\cite{BR,HH}.

\begin{theorem}[Cubic multiplicity bounds]\label{cubic-bounds}
Let $R$ be a three-dimensional Cohen--Macaulay normal local domain and
$R\subseteq S$ a torsion-free module-finite extension of rank three
with $\pd_R S\le1$. Put $E=\Ext^1_R(S/R,R)$,
$e=e_{\m}(R)$ and $v=\dim_k\m/\m^2$.
If $\ell_R(E)<4e$, or if $R$ is singular Gorenstein and
$v\ell_R(\m^2E)<4(e-1)$, the extension splits and $\mu_R(S)\le4$.
\end{theorem}
\begin{proof}
Put $C=S/R$.  This module is torsion-free.  Indeed, if $as\in R$
for $0\ne a\in R$ and $s\in S$, then in $S\otimes_R\operatorname{Frac}R$
the element $s$ is a scalar.  It is integral over $R$, hence belongs
to $R$ by normality.  The sequence
$0\to R\to S\to C\to0$ gives $\pd_R C\leq1$.

Under either numerical hypothesis
$E=\Ext^1_R(C,R)$ has finite length.  For the second hypothesis,
this follows because $\m^2E$ has finite length by assumption and
$E/\m^2E$ has finite length automatically.  Localizing a free
resolution of $C$ away from $\m$, the vanishing of $E$ makes the
dual presentation surjective and hence split.  Thus $C$ is free
on the punctured spectrum.  At the maximal ideal its depth is at
least two by Auslander--Buchsbaum.  Therefore $C$ is reflexive,
of rank two, and its determinant is trivial by finite projective
dimension.  Exterior product, followed by a choice of this
trivialization, gives
\begin{equation}\label{cubic-quotient-self-dual}
 C\simeq C^*.
\end{equation}
To justify this last isomorphism, the exterior pairing is perfect
at height-one primes and both modules are reflexive. The resulting
codimension-one isomorphism extends over $R$.

Let $b=\dim_k(E/\m E)$.  If $b>0$, the minimal resolution of $C$
and its dual have the forms
\begin{equation}\label{cubic-quotient-resolutions}
 \begin{gathered}
 0\longrightarrow R^b\longrightarrow R^{b+2}
       \longrightarrow C\longrightarrow0,\\
 0\longrightarrow C^*\longrightarrow R^{b+2}
       \longrightarrow R^b\longrightarrow E\longrightarrow0.
 \end{gathered}
\end{equation}
Let $N$ be the image in $R^b$.  It has finite colength and is
contained in $\m R^b$.  If $q=\mu_R(N)$, then $q\leq b+2$.
On the other hand, the ideal of maximal minors of a minimal
$b\times q$ matrix for $N$ is $\m$-primary, because $R^b/N=E$
is nonzero and of finite length.  The determinantal height bound
therefore gives $3\leq q-b+1$, so $q=b+2$.
This is exactly the number required for $N$ to be a parameter
module in a free module of rank $b$ over a three-dimensional base.

For such a parameter module over a Cohen--Macaulay ring, the
Buchsbaum--Rim formula gives
$\ell_R(R^b/N)=e_{\mathrm{BR}}(N,R^b)$ \cite{BR,HH}.
Here $e_{\mathrm{BR}}$ is the multiplicity determined by the leading
coefficient of
\[
 j\longmapsto
 \ell_R\bigl(\operatorname{Sym}^j(R^b)/N^j\bigr),
\]
where $N^j$ is the degree-$j$ image of the symmetric algebra of
$N$ in $\operatorname{Sym}(R^b)$: its leading term is
$e_{\mathrm{BR}}j^{b+2}/(b+2)!$.
The inclusion $N\subseteq\m R^b$ compares these colengths in every
degree and yields
\begin{equation}\label{cubic-BR-bound}
 \ell_R(E)=e_{\mathrm{BR}}(N,R^b)
     \geq e_{\mathrm{BR}}(\m R^b,R^b)
     =e\binom{b+2}{3}.
\end{equation}
For the last equality, the degree-$j$ colength for $\m R^b$ is
\[
 \binom{j+b-1}{b-1}\ell_R(R/\m^j).
\]
Since $\ell_R(R/\m^j)$ has leading term $ej^3/3!$, comparison
of the leading coefficients proves the equality.  As
$\binom{4}{3}=4$, the inequality $\ell_R(E)<4e$ forces $b\leq1$.

For the second criterion assume that $R$ is Gorenstein.
Two dimension shifts in \eqref{cubic-quotient-resolutions}, followed
by \eqref{cubic-quotient-self-dual}, give
\[
 \Ext^3_R(E,R)\simeq\Ext^1_R(C^*,R)\simeq E.
\]
Local duality identifies the first module with the Matlis dual
of $E$.  Completion preserves all the finite
lengths and the displayed isomorphisms.
Thus $E$ is Matlis self-dual, and its socle
$(0:_E\m)$ has length $\dim_k(E/\m E)=b$.
If $x_1,\ldots,x_v$ minimally generate $\m$, the map
\[
 \m E\longrightarrow(\m^2E)^v,
 \qquad u\longmapsto(x_1u,\ldots,x_vu),
\]
has kernel contained in that socle.  It follows that
\begin{equation}\label{cubic-socle-bound}
 \ell_R(E)\leq2b+v\ell_R(\m^2E).
\end{equation}
For $b\geq2$ one has $\binom{b+2}{3}\geq2b$.
Combining \eqref{cubic-BR-bound} and \eqref{cubic-socle-bound}
therefore gives
\[
 v\ell_R(\m^2E)\geq2b(e-1)\geq4(e-1).
\]
The last inequality uses $e\geq2$, since the Cohen--Macaulay base
is singular.  This contradicts the second hypothesis, so again
$b\leq1$.

If $b=0$, the quotient $C$ is free and the exact sequence
$0\to R\to S\to C\to0$ splits.  Suppose $b=1$.
Choose lifts $s_1,s_2,s_3\in S$ of the three minimal generators
of $C$.  If $(d_1,d_2,d_3)^t$ generates its relation module,
then $\sum_{i=1}^3d_i s_i\in R$.  Write this scalar as $-d_0$.
The generators $1_S,s_1,s_2,s_3$ present $S$ with the single
relation
\[
 d=(d_0,d_1,d_2,d_3)^t.
\]
Indeed any relation projects to a multiple of the given relation
of $C$, and subtraction leaves a scalar relation, which is zero
by injectivity of $R\subseteq S$.

Lift multiplication by $s_1$ to a matrix $P$ on this free cover,
choosing its first column to be the basis vector corresponding
to $s_1$.  Preservation of the relation module gives $Pd=\lambda d$
for some $\lambda\in R$.  Its row indexed by $s_1$ says
\[
 d_0+P_{11}d_1+P_{12}d_2+P_{13}d_3=\lambda d_1.
\]
Consequently $d_0\in(d_1,d_2,d_3)$.  A row vector with first
entry $1$ can therefore be chosen to kill $d$ and descends to
an $R$-linear retraction of the unit.  This last step is where
the algebra structure, beyond the rank-two module $C$, is used.
Finally $\mu_R(S)\leq1+\mu_R(C)=b+3\leq4$ under either bound.
\end{proof}
